\chapter{浏览器控制}

\section{浏览器选项配置}

在进行Web自动化测试时，合理配置浏览器选项可以显著提高脚本的稳定性和执行效率。不同浏览器提供了丰富的配置选项来满足各种测试需求。

\subsection{Chrome浏览器配置}

\begin{lstlisting}[style=demostyle,numbers=none]
from selenium import webdriver
from selenium.webdriver.chrome.options import Options
import os

# 创建Chrome选项对象
chrome_options = Options()

# 无头模式（后台运行，不显示浏览器界面）
chrome_options.add_argument("--headless")

# 窗口大小设置
chrome_options.add_argument("--window-size=1920,1080")

# 禁用GPU加速（在某些环境下可以避免问题）
chrome_options.add_argument("--disable-gpu")

# 禁用开发者工具
chrome_options.add_argument("--disable-dev-shm-usage")

# 禁用沙盒模式（在Docker等环境中可能需要）
chrome_options.add_argument("--no-sandbox")

# 设置用户代理
chrome_options.add_argument("--user-agent=Mozilla/5.0 (Windows NT 10.0; Win64; x64) AppleWebKit/537.36")

# 禁用图片加载（提高页面加载速度）
prefs = {"profile.managed_default_content_settings.images": 2}
chrome_options.add_experimental_option("prefs", prefs)

# 禁用通知
chrome_options.add_argument("--disable-notifications")

# 禁用弹出窗口阻止程序
chrome_options.add_argument("--disable-popup-blocking")

# 设置下载目录
download_dir = os.path.join(os.getcwd(), "downloads")
prefs = {
    "download.default_directory": download_dir,
    "download.prompt_for_download": False,
    "download.directory_upgrade": True
}
chrome_options.add_experimental_option("prefs", prefs)

# 启动最大化窗口
chrome_options.add_argument("--start-maximized")

# 创建浏览器实例
driver = webdriver.Chrome(options=chrome_options)
\end{lstlisting}

\subsection{Firefox浏览器配置}

\begin{lstlisting}[style=demostyle,numbers=none]
from selenium import webdriver
from selenium.webdriver.firefox.options import Options
from selenium.webdriver.firefox.firefox_profile import FirefoxProfile

# Firefox选项配置
firefox_options = Options()

# 无头模式
firefox_options.add_argument("--headless")

# 窗口大小
firefox_options.add_argument("--width=1920")
firefox_options.add_argument("--height=1080")

# 创建Firefox配置文件
profile = FirefoxProfile()

# 禁用图片加载
profile.set_preference("permissions.default.image", 2)

# 设置下载行为
profile.set_preference("browser.download.folderList", 2)
profile.set_preference("browser.download.dir", "/path/to/download")
profile.set_preference("browser.helperApps.neverAsk.saveToDisk", 
                      "application/pdf,application/octet-stream")

# 禁用Firefox更新提示
profile.set_preference("app.update.enabled", False)

# 设置用户代理
profile.set_preference("general.useragent.override", 
                      "Mozilla/5.0 (Windows NT 10.0; Win64; x64; rv:91.0)")

# 应用配置
firefox_options.profile = profile

driver = webdriver.Firefox(options=firefox_options)
\end{lstlisting}

\subsection{Edge浏览器配置}

\begin{lstlisting}[style=demostyle,numbers=none]
from selenium import webdriver
from selenium.webdriver.edge.options import Options

# Edge选项配置
edge_options = Options()

# 基础配置
edge_options.add_argument("--headless")
edge_options.add_argument("--window-size=1920,1080")
edge_options.add_argument("--disable-gpu")

# Edge特有配置
edge_options.add_argument("--disable-extensions")
edge_options.add_argument("--disable-plugins")

# 实验性选项
edge_options.add_experimental_option("useAutomationExtension", False)
edge_options.add_experimental_option("excludeSwitches", ["enable-automation"])

driver = webdriver.Edge(options=edge_options)
\end{lstlisting}

\section{多窗口和标签页管理}

现代Web应用经常涉及多个窗口和标签页的操作，掌握这些技能对于复杂场景的自动化测试至关重要。

\subsection{窗口句柄操作}

\begin{lstlisting}[style=demostyle,numbers=none]
from selenium import webdriver
from selenium.webdriver.common.by import By
import time

driver = webdriver.Chrome()
driver.get("https://example.com")

# 获取当前窗口句柄
main_window = driver.current_window_handle
print(f"主窗口句柄: {main_window}")

# 获取所有窗口句柄
all_windows = driver.window_handles
print(f"所有窗口句柄: {all_windows}")

# 点击链接打开新窗口
new_window_link = driver.find_element(By.LINK_TEXT, "打开新窗口")
new_window_link.click()

# 等待新窗口出现
time.sleep(2)

# 切换到新窗口
for window_handle in driver.window_handles:
    if window_handle != main_window:
        driver.switch_to.window(window_handle)
        break

print(f"当前窗口标题: {driver.title}")

# 在新窗口中执行操作
# ... 执行相关操作 ...

# 关闭当前窗口
driver.close()

# 切换回主窗口
driver.switch_to.window(main_window)

# 验证已回到主窗口
print(f"回到主窗口标题: {driver.title}")
\end{lstlisting}

\subsection{标签页管理}

\begin{lstlisting}[style=demostyle,numbers=none]
# 打开新标签页的几种方法

# 方法1：使用JavaScript打开新标签页
driver.execute_script("window.open('https://www.google.com', '_blank');")

# 方法2：使用键盘快捷键
from selenium.webdriver.common.keys import Keys
from selenium.webdriver.common.action_chains import ActionChains

# Ctrl+T 打开新标签页
ActionChains(driver).key_down(Keys.CONTROL).send_keys('t').key_up(Keys.CONTROL).perform()

# 方法3：使用Selenium 4的新方法
driver.switch_to.new_window('tab')  # 打开新标签页
driver.switch_to.new_window('window')  # 打开新窗口

# 标签页管理函数
class TabManager:
    def __init__(self, driver):
        self.driver = driver
        self.tabs = {}
    
    def open_new_tab(self, url, name):
        """打开新标签页并命名"""
        self.driver.execute_script(f"window.open('{url}', '_blank');")
        self.driver.switch_to.window(self.driver.window_handles[-1])
        self.tabs[name] = self.driver.current_window_handle
        return self.driver.current_window_handle
    
    def switch_to_tab(self, name):
        """切换到指定标签页"""
        if name in self.tabs:
            self.driver.switch_to.window(self.tabs[name])
        else:
            raise ValueError(f"标签页 '{name}' 不存在")
    
    def close_tab(self, name):
        """关闭指定标签页"""
        if name in self.tabs:
            self.driver.switch_to.window(self.tabs[name])
            self.driver.close()
            del self.tabs[name]
    
    def close_all_tabs_except(self, keep_tab):
        """关闭除指定标签页外的所有标签页"""
        for name, handle in self.tabs.items():
            if name != keep_tab:
                self.driver.switch_to.window(handle)
                self.driver.close()
        self.tabs = {keep_tab: self.tabs[keep_tab]}

# 使用标签页管理器
tab_manager = TabManager(driver)
tab_manager.open_new_tab("https://www.baidu.com", "baidu")
tab_manager.open_new_tab("https://www.google.com", "google")

tab_manager.switch_to_tab("baidu")
print(f"当前页面标题: {driver.title}")

tab_manager.switch_to_tab("google")
print(f"当前页面标题: {driver.title}")
\end{lstlisting}

\subsection{窗口位置和大小控制}

\begin{lstlisting}[style=demostyle,numbers=none]
# 窗口大小和位置操作

# 设置窗口大小
driver.set_window_size(1200, 800)

# 获取窗口大小
size = driver.get_window_size()
print(f"窗口大小: 宽度={size['width']}, 高度={size['height']}")

# 设置窗口位置
driver.set_window_position(100, 100)

# 获取窗口位置
position = driver.get_window_position()
print(f"窗口位置: X={position['x']}, Y={position['y']}")

# 最大化窗口
driver.maximize_window()

# 最小化窗口
driver.minimize_window()

# 全屏模式
driver.fullscreen_window()

# 获取窗口矩形（位置和大小）
rect = driver.get_window_rect()
print(f"窗口矩形: {rect}")

# 设置窗口矩形
driver.set_window_rect(x=50, y=50, width=1400, height=900)

# 多显示器支持
def move_to_second_monitor(driver):
    """将窗口移动到第二个显示器"""
    # 假设第二个显示器在右侧，分辨率为1920x1080
    driver.set_window_position(1920, 0)
    driver.set_window_size(1200, 800)
\end{lstlisting}

\section{页面导航}

页面导航是Web自动化中的基础操作，包括页面跳转、历史记录操作和页面刷新等。

\subsection{基础导航操作}

\begin{lstlisting}[style=demostyle,numbers=none]
# 导航到指定URL
driver.get("https://www.example.com")

# 获取当前URL
current_url = driver.current_url
print(f"当前URL: {current_url}")

# 获取页面标题
title = driver.title
print(f"页面标题: {title}")

# 后退
driver.back()

# 前进
driver.forward()

# 刷新页面
driver.refresh()

# 获取页面源代码
page_source = driver.page_source
print(f"页面源代码长度: {len(page_source)}")
\end{lstlisting}

\subsection{等待页面加载}

\begin{lstlisting}[style=demostyle,numbers=none]
from selenium.webdriver.support.ui import WebDriverWait
from selenium.webdriver.support import expected_conditions as EC
from selenium.webdriver.common.by import By

# 等待页面标题包含特定文本
wait = WebDriverWait(driver, 10)
wait.until(EC.title_contains("预期标题"))

# 等待URL变化
wait.until(EC.url_contains("/dashboard"))

# 等待特定元素出现
wait.until(EC.presence_of_element_located((By.ID, "content")))

# 自定义页面加载完成检查
def page_loaded(driver):
    """检查页面是否完全加载"""
    return driver.execute_script("return document.readyState") == "complete"

# 等待页面完全加载
wait.until(page_loaded)

# 等待Ajax请求完成（jQuery）
def ajax_complete(driver):
    """等待jQuery Ajax请求完成"""
    try:
        return driver.execute_script("return jQuery.active == 0")
    except:
        return True

wait.until(ajax_complete)
\end{lstlisting}

\subsection{页面滚动操作}

\begin{lstlisting}[style=demostyle,numbers=none]
# 滚动到页面底部
driver.execute_script("window.scrollTo(0, document.body.scrollHeight);")

# 滚动到页面顶部
driver.execute_script("window.scrollTo(0, 0);")

# 滚动到指定位置
driver.execute_script("window.scrollTo(0, 500);")

# 滚动到指定元素
element = driver.find_element(By.ID, "target-element")
driver.execute_script("arguments[0].scrollIntoView();", element)

# 平滑滚动到元素
driver.execute_script("arguments[0].scrollIntoView({behavior: 'smooth'});", element)

# 水平滚动
driver.execute_script("window.scrollBy(500, 0);")

# 在可滚动容器内滚动
scrollable_div = driver.find_element(By.ID, "scrollable-content")
driver.execute_script("arguments[0].scrollTop = arguments[0].scrollHeight;", scrollable_div)

# 分步滚动（模拟用户滚动行为）
def smooth_scroll_to_bottom(driver, step=300, delay=0.5):
    """平滑滚动到页面底部"""
    import time
    
    last_height = driver.execute_script("return document.body.scrollHeight")
    current_position = 0
    
    while current_position < last_height:
        current_position += step
        driver.execute_script(f"window.scrollTo(0, {current_position});")
        time.sleep(delay)
        
        # 检查页面高度是否发生变化（动态加载内容）
        new_height = driver.execute_script("return document.body.scrollHeight")
        if new_height > last_height:
            last_height = new_height

# 使用ActionChains进行滚动
from selenium.webdriver.common.action_chains import ActionChains

element = driver.find_element(By.ID, "content")
actions = ActionChains(driver)
actions.scroll_to_element(element).perform()

# 鼠标滚轮滚动
actions.scroll_by_amount(0, 1000).perform()  # 垂直滚动1000像素
\end{lstlisting}

\section{JavaScript代码执行}

Selenium提供了强大的JavaScript执行能力，可以直接在浏览器中运行JavaScript代码，实现复杂的操作。

\subsection{基础JavaScript执行}

\begin{lstlisting}[style=demostyle,numbers=none]
# 执行简单JavaScript代码
result = driver.execute_script("return document.title;")
print(f"页面标题: {result}")

# 执行无返回值的JavaScript
driver.execute_script("alert('Hello from Selenium!');")

# 传递参数给JavaScript
element = driver.find_element(By.ID, "username")
driver.execute_script("arguments[0].style.border = '3px solid red';", element)

# 多参数传递
name = "测试用户"
age = 25
driver.execute_script("arguments[0].value = arguments[1] + ' - ' + arguments[2];", 
                     element, name, age)

# 返回复杂数据类型
page_info = driver.execute_script("""
    return {
        title: document.title,
        url: window.location.href,
        width: window.innerWidth,
        height: window.innerHeight,
        scrollY: window.scrollY
    };
""")
print(f"页面信息: {page_info}")
\end{lstlisting}

\subsection{DOM操作}

\begin{lstlisting}[style=demostyle,numbers=none]
# 修改元素内容
driver.execute_script("document.getElementById('demo').innerHTML = '新内容';")

# 修改元素样式
element = driver.find_element(By.ID, "target")
driver.execute_script("""
    arguments[0].style.backgroundColor = 'yellow';
    arguments[0].style.fontSize = '20px';
    arguments[0].style.fontWeight = 'bold';
""", element)

# 添加新元素
driver.execute_script("""
    var newDiv = document.createElement('div');
    newDiv.innerHTML = '这是动态添加的内容';
    newDiv.id = 'dynamic-content';
    document.body.appendChild(newDiv);
""")

# 移除元素
driver.execute_script("document.getElementById('unwanted-element').remove();")

# 获取元素属性
attributes = driver.execute_script("""
    var elem = arguments[0];
    return {
        id: elem.id,
        className: elem.className,
        tagName: elem.tagName,
        innerText: elem.innerText,
        offsetWidth: elem.offsetWidth,
        offsetHeight: elem.offsetHeight
    };
""", element)
print(f"元素属性: {attributes}")

# 触发事件
driver.execute_script("arguments[0].click();", element)
driver.execute_script("arguments[0].focus();", element)

# 自定义事件
driver.execute_script("""
    var event = new Event('custom-event');
    arguments[0].dispatchEvent(event);
""", element)
\end{lstlisting}

\subsection{页面信息获取}

\begin{lstlisting}[style=demostyle,numbers=none]
# 获取页面加载状态
ready_state = driver.execute_script("return document.readyState;")
print(f"页面状态: {ready_state}")

# 获取视口信息
viewport = driver.execute_script("""
    return {
        width: window.innerWidth,
        height: window.innerHeight,
        scrollX: window.scrollX,
        scrollY: window.scrollY,
        devicePixelRatio: window.devicePixelRatio
    };
""")

# 获取所有链接
links = driver.execute_script("""
    var links = document.getElementsByTagName('a');
    var linkData = [];
    for(var i = 0; i < links.length; i++) {
        linkData.push({
            text: links[i].innerText,
            href: links[i].href,
            target: links[i].target
        });
    }
    return linkData;
""")

# 获取表单数据
form_data = driver.execute_script("""
    var forms = document.forms;
    var formData = [];
    for(var i = 0; i < forms.length; i++) {
        var form = forms[i];
        var inputs = form.querySelectorAll('input, select, textarea');
        var fields = [];
        for(var j = 0; j < inputs.length; j++) {
            fields.push({
                name: inputs[j].name,
                type: inputs[j].type,
                value: inputs[j].value
            });
        }
        formData.push({
            action: form.action,
            method: form.method,
            fields: fields
        });
    }
    return formData;
""")

# 检查元素可见性
is_visible = driver.execute_script("""
    var elem = arguments[0];
    var rect = elem.getBoundingClientRect();
    return rect.width > 0 && rect.height > 0 && 
           rect.top >= 0 && rect.left >= 0 &&
           rect.bottom <= window.innerHeight &&
           rect.right <= window.innerWidth;
""", element)
\end{lstlisting}

\subsection{异步JavaScript执行}

\begin{lstlisting}[style=demostyle,numbers=none]
# 执行异步JavaScript（已在Selenium 4中移除，但可以用其他方式实现）
# 使用Promise来处理异步操作

# 等待Ajax请求完成
def wait_for_ajax_complete(driver, timeout=10):
    """等待所有Ajax请求完成"""
    return driver.execute_script("""
        return new Promise((resolve) => {
            if (typeof jQuery !== 'undefined') {
                let checkAjax = () => {
                    if (jQuery.active === 0) {
                        resolve(true);
                    } else {
                        setTimeout(checkAjax, 100);
                    }
                };
                checkAjax();
            } else {
                resolve(true);
            }
        });
    """)

# 等待自定义条件
def wait_for_condition(driver, condition_script, timeout=10):
    """等待自定义JavaScript条件为真"""
    import time
    start_time = time.time()
    
    while time.time() - start_time < timeout:
        try:
            result = driver.execute_script(f"return ({condition_script});")
            if result:
                return True
        except:
            pass
        time.sleep(0.1)
    
    return False

# 使用示例
condition = "document.querySelectorAll('.dynamic-content').length > 0"
if wait_for_condition(driver, condition):
    print("动态内容已加载")
\end{lstlisting}

\section{Cookie和Session管理}

Cookie和Session管理对于处理用户认证、保持登录状态和测试不同用户场景非常重要。

\subsection{Cookie操作}

\begin{lstlisting}[style=demostyle,numbers=none]
# 获取所有Cookie
all_cookies = driver.get_cookies()
print(f"所有Cookie: {all_cookies}")

# 获取特定Cookie
session_cookie = driver.get_cookie("sessionid")
if session_cookie:
    print(f"Session Cookie: {session_cookie}")

# 添加Cookie
driver.add_cookie({
    "name": "test_cookie",
    "value": "test_value",
    "domain": "example.com",
    "path": "/",
    "httpOnly": False,
    "secure": False
})

# 添加带有过期时间的Cookie
import time
from datetime import datetime, timedelta

expire_time = int((datetime.now() + timedelta(days=7)).timestamp())
driver.add_cookie({
    "name": "persistent_cookie",
    "value": "persistent_value",
    "expiry": expire_time
})

# 删除特定Cookie
driver.delete_cookie("test_cookie")

# 删除所有Cookie
driver.delete_all_cookies()

# Cookie管理类
class CookieManager:
    def __init__(self, driver):
        self.driver = driver
    
    def save_cookies(self, filepath):
        """保存Cookie到文件"""
        import json
        cookies = self.driver.get_cookies()
        with open(filepath, 'w') as f:
            json.dump(cookies, f)
    
    def load_cookies(self, filepath):
        """从文件加载Cookie"""
        import json
        try:
            with open(filepath, 'r') as f:
                cookies = json.load(f)
            
            for cookie in cookies:
                # 移除可能导致问题的字段
                if 'expiry' in cookie:
                    cookie['expiry'] = int(cookie['expiry'])
                self.driver.add_cookie(cookie)
        except FileNotFoundError:
            print(f"Cookie文件 {filepath} 不存在")
    
    def get_cookie_dict(self):
        """获取Cookie字典格式"""
        cookies = self.driver.get_cookies()
        return {cookie['name']: cookie['value'] for cookie in cookies}
    
    def set_login_cookies(self, user_id, session_token):
        """设置登录相关Cookie"""
        self.driver.add_cookie({"name": "user_id", "value": str(user_id)})
        self.driver.add_cookie({"name": "session_token", "value": session_token})
        self.driver.add_cookie({"name": "logged_in", "value": "true"})

# 使用Cookie管理器
cookie_manager = CookieManager(driver)

# 登录后保存Cookie
driver.get("https://example.com/login")
# ... 执行登录操作 ...
cookie_manager.save_cookies("login_cookies.json")

# 下次直接加载Cookie实现免登录
driver.get("https://example.com")
cookie_manager.load_cookies("login_cookies.json")
driver.refresh()  # 刷新页面使Cookie生效
\end{lstlisting}

\subsection{Session管理}

\begin{lstlisting}[style=demostyle,numbers=none]
# Session状态检查
def check_login_status(driver):
    """检查登录状态"""
    try:
        # 方法1：检查特定元素是否存在
        user_menu = driver.find_element(By.CLASS_NAME, "user-menu")
        return True
    except:
        try:
            # 方法2：检查Cookie
            session_cookie = driver.get_cookie("session_id")
            return session_cookie is not None
        except:
            return False

# 自动登录函数
def auto_login(driver, username, password):
    """自动登录函数"""
    driver.get("https://example.com/login")
    
    # 输入用户名和密码
    username_field = driver.find_element(By.NAME, "username")
    password_field = driver.find_element(By.NAME, "password")
    
    username_field.send_keys(username)
    password_field.send_keys(password)
    
    # 点击登录按钮
    login_button = driver.find_element(By.ID, "login-button")
    login_button.click()
    
    # 等待登录完成
    from selenium.webdriver.support.ui import WebDriverWait
    from selenium.webdriver.support import expected_conditions as EC
    
    wait = WebDriverWait(driver, 10)
    try:
        # 等待跳转到用户页面或出现用户菜单
        wait.until(EC.any_of(
            EC.url_contains("/dashboard"),
            EC.presence_of_element_located((By.CLASS_NAME, "user-menu"))
        ))
        return True
    except:
        return False

# 多用户Session管理
class MultiUserSessionManager:
    def __init__(self):
        self.sessions = {}
    
    def create_session(self, user_id, driver_options=None):
        """为用户创建新的浏览器会话"""
        if driver_options is None:
            driver_options = webdriver.ChromeOptions()
            driver_options.add_argument(f"--user-data-dir=./user_data_{user_id}")
        
        driver = webdriver.Chrome(options=driver_options)
        self.sessions[user_id] = driver
        return driver
    
    def switch_to_user(self, user_id):
        """切换到指定用户的会话"""
        return self.sessions.get(user_id)
    
    def login_user(self, user_id, username, password):
        """为指定用户执行登录"""
        driver = self.sessions.get(user_id)
        if driver:
            return auto_login(driver, username, password)
        return False
    
    def close_session(self, user_id):
        """关闭用户会话"""
        if user_id in self.sessions:
            self.sessions[user_id].quit()
            del self.sessions[user_id]
    
    def close_all_sessions(self):
        """关闭所有会话"""
        for driver in self.sessions.values():
            driver.quit()
        self.sessions.clear()

# 使用多用户管理器
session_manager = MultiUserSessionManager()

# 创建多个用户会话
admin_driver = session_manager.create_session("admin")
user_driver = session_manager.create_session("user1")

# 分别登录不同用户
session_manager.login_user("admin", "admin@example.com", "admin_password")
session_manager.login_user("user1", "user1@example.com", "user_password")

# 在不同会话中执行操作
admin_driver.get("https://example.com/admin")
user_driver.get("https://example.com/profile")
\end{lstlisting}

\subsection{高级会话管理}

\begin{lstlisting}[style=demostyle,numbers=none]
# 会话持久化
class PersistentSession:
    def __init__(self, session_file):
        self.session_file = session_file
        self.driver = None
    
    def start_session(self, restore=True):
        """启动会话，可选择是否恢复之前的状态"""
        options = webdriver.ChromeOptions()
        
        if restore:
            # 使用固定的用户数据目录
            options.add_argument(f"--user-data-dir=./session_data")
        
        self.driver = webdriver.Chrome(options=options)
        
        if restore:
            self.restore_session()
        
        return self.driver
    
    def save_session(self):
        """保存当前会话状态"""
        if not self.driver:
            return
        
        session_data = {
            'url': self.driver.current_url,
            'cookies': self.driver.get_cookies(),
            'local_storage': self.driver.execute_script("return localStorage;"),
            'session_storage': self.driver.execute_script("return sessionStorage;")
        }
        
        import json
        with open(self.session_file, 'w') as f:
            json.dump(session_data, f)
    
    def restore_session(self):
        """恢复会话状态"""
        try:
            import json
            with open(self.session_file, 'r') as f:
                session_data = json.load(f)
            
            # 导航到之前的URL
            self.driver.get(session_data['url'])
            
            # 恢复Cookie
            for cookie in session_data['cookies']:
                try:
                    self.driver.add_cookie(cookie)
                except:
                    pass
            
            # 恢复LocalStorage
            for key, value in session_data.get('local_storage', {}).items():
                self.driver.execute_script(
                    f"localStorage.setItem('{key}', '{value}');"
                )
            
            # 恢复SessionStorage
            for key, value in session_data.get('session_storage', {}).items():
                self.driver.execute_script(
                    f"sessionStorage.setItem('{key}', '{value}');"
                )
            
            # 刷新页面使状态生效
            self.driver.refresh()
            
        except FileNotFoundError:
            print("未找到会话文件，将创建新会话")
    
    def end_session(self, save=True):
        """结束会话"""
        if save and self.driver:
            self.save_session()
        
        if self.driver:
            self.driver.quit()

# 使用持久化会话
persistent_session = PersistentSession("my_session.json")
driver = persistent_session.start_session(restore=True)

# 执行操作...
driver.get("https://example.com")

# 结束时保存会话
persistent_session.end_session(save=True)
\end{lstlisting}